\section{User Study}
\label{sec:8}

The closed-loop stage will test whether decisions improve interaction, not merely whether support-need labels can be predicted. Proactive timing research motivates comparing alternative intervention timing, and intervention work motivates measuring consequences for the ongoing task as well as physiological response ~\cite{sensing,calm}. The study design below is a proposed evaluation logic; allocation, power, parameters, and findings remain to be supplied.

\subsection{Conditions and isolation of mechanisms}

The proposed core conditions are no agent support, an agent using behavioral and task information, and an agent additionally using multi-person physiology. The two agents will share the same available actions, interface, and prompt library. A timing control with comparable intervention content and frequency but state-unrelated eligible timing will be included if feasible. This comparison is motivated by prior work evaluating aligned and alternative timing ~\cite{sensing}; the control will respect task constraints rather than interrupt at arbitrary unsafe moments.

\missing{[TO COMPLETE: conditions, team-level allocation, between-team or counterbalanced design, carryover controls, matched task difficulty, prompt matching, primary contrast, and manipulation checks.]}

\subsection{Participants and procedure}

\missing{[TO COMPLETE: recruitment and eligibility, team composition and familiarity, independent team sample size, pilot-based power or precision analysis, compensation, ethics approval, familiarization, task order, intervention exposure, and interview procedure.]}

\subsection{Outcomes}

For the primary collaboration study, the proposed primary outcome is successful sharing and uptake of task-relevant information within a prespecified interval. Secondary outcomes will assess decision quality, unwanted or declined prompts, interruption cost, perceived agency, workload, and participant explanations of helpful and unhelpful episodes. These choices draw on collaborative measurement, moderation, and proactive-interaction research ~\cite{p57,p70,expect}. Physiology will remain an explanatory measure rather than the definition of successful collaboration ~\cite{p32}.

\missing{[TO COMPLETE: primary outcome and denominator, task scoring rubric, communication coding, validated questionnaire choices, prompt-appropriateness assessment, coder blinding, and interview guide.]}

\subsection{Analysis}

The proposed analysis will treat teams as the independent allocation unit and model repeated sessions or events and members within teams as appropriate. Effect estimates and uncertainty will accompany hypothesis tests. Event-level differences will not automatically be interpreted causally because earlier interventions can change which later events occur. The primary randomized-condition comparison will therefore be distinguished from exploratory analyses of the mechanism.

\missing{[TO COMPLETE: estimands, statistical models and distributions, random effects, covariates, missing-data treatment, multiplicity handling, equivalence or precision criteria for null findings, and qualitative coding procedure.]}

\subsection{Closed loop results}

\missing{[RESULTS TO COMPLETE: sample flow; primary condition contrast with effect size and uncertainty; appropriate and unwanted prompts; task outcomes; user agency; mechanism checks; qualitative themes and counterexamples. Include negative and mixed findings.]}
